\documentclass{article}
\usepackage{amsmath,amssymb,amsthm}
\usepackage{url}

\newtheorem{definition}{Definition}
\newtheorem{lemma}[definition]{Lemma}
\newcommand{\source}[1]{\par\noindent\textbf{Source.} #1}
\newcommand{\lean}[1]{\par\noindent\textbf{Lean target.} \texttt{#1}}
\newcommand{\status}[1]{\par\noindent\textbf{Status.} \texttt{\detokenize{#1}}}

\begin{document}

\section*{Blueprint 002: Canonical Stabilizer State}

\subsection*{Scope}

This note supports target `STATE-001'.  It formalizes only the Pauli-label
coefficient rule for the canonical stabilizer projection state attached to an
isotropic stabilizer label subspace and a consistent sign character.  It does
not yet formalize the Pauli \(C^*\)-algebra, positivity of the resulting
linear functional, GNS data, or any von Neumann algebra.

\subsection*{Conventions}

Let \(S=(\Gamma,W)\) be the algebraic stabilizer system of `SYS-001', so
\[
  W\leq V_\Gamma
  =
  \mathbb F_2^{(\Gamma)}\oplus \mathbb F_2^{(\Gamma)}
\]
is isotropic.  A stabilizer character is a group homomorphism
\[
  \chi:W\to\{\pm 1\}
\]
from the additive stabilizer label group to the multiplicative sign group.
Equivalently, in Lean we may encode the same character as an additive map
\[
  \bar\chi:W\to\mathbb F_2
\]
and read the sign as \( (-1)^{\bar\chi(w)} \).

\begin{definition}[CanonicalStatePackage, ID: STATE-001, project proof note 004, stabilizer character, crossing-stabilizer locality, and product form; Gottesman Chapter 3.5, Definition 3.13, PDF p.~56]
Given an isotropic stabilizer label subspace \(W\leq V_\Gamma\) and a
stabilizer character \(\chi:W\to\{\pm1\}\), the canonical coefficient state on
Pauli labels is determined on basis Pauli labels by
\[
  \omega(P_v)=
  \begin{cases}
    \chi(v), & v\in W,\\
    0, & v\notin W.
  \end{cases}
\]
\source{Project proof note
\path{../runs/2026-06-18-stabilizer-algebra-classification/artifacts/004-internal-stabilizer-reduction-type-classification-R1.tex},
``Stabilizer character'', ``Crossing stabilizers do not change the local
coefficients'', and ``Product form after symplectic basis reduction'', for the
canonical projection-state Pauli coefficient rule and its character convention.
Gottesman, \emph{Surviving as a Quantum Computer in a Classical World},
Chapter 3.5, Definition 3.13, PDF p.~56, for stabilizer labels as commuting
binary symplectic subspaces.}
\lean{Formalization.Stabilizer.CanonicalStatePackage}
\status{blueprint_note_complete}
\end{definition}

\begin{proof}[Informal well-formedness and state-origin check]
The formula is well formed at the Pauli-label coefficient level because every
label \(v\in V_\Gamma\) either lies in the stabilizer subspace \(W\) or does
not.  In the first case the coefficient is the sign prescribed by the
character \(\chi\); in the second case the coefficient is zero.  The character
condition is the consistency condition for assigning simultaneous stabilizer
eigenvalues to commuting stabilizer labels: if \(w,w'\in W\), then the
coefficient assigned to the product label \(w+w'\) is
\(\chi(w+w')=\chi(w)\chi(w')\).  The isotropy of \(W\) from `SYS-001' is the
binary-symplectic commutation condition needed for this simultaneous
assignment.

The project proof note verifies the operator-algebraic origin of this
coefficient rule for the canonical projection state and shows that, after an
adapted symplectic basis reduction, the restriction is a tensor product of
one-site \(Z\)-eigenstates on fixed modes and normalized traces on unfixed
modes.  That product-form argument is the positivity/state justification for
the informal mathematics.  The present Lean target records only the
blueprint-backed algebraic coefficient package; formal positivity as a
\(C^*\)-state is reserved for the later Pauli-algebra and GNS targets.
\end{proof}

\end{document}
